\chapter{Smelly Feet (Bromodosis)}

\section*{Definition}
Smelly feet, medically called bromodosis, is an unpleasant foot odor caused by moisture and bacteria interacting with sweat on the skin. It is a very common condition that usually improves with good foot hygiene.

\section*{What Happens in Your Body}
Your feet have more sweat glands per square inch than any other part of your body — about 250,000 sweat glands per foot. Sweat itself is odorless, but when your feet are enclosed in shoes, the moisture creates a warm, dark environment where bacteria and fungi thrive. These microorganisms break down dead skin cells and components of sweat, producing waste products that release the characteristic unpleasant smell. The bacteria that cause the odor are typically species of Brevibacterium, Staphylococcus, and Propionibacterium.

\section*{Causes}
\begin{itemize}
  \item Excessive sweating (hyperhidrosis) of the feet — either genetic or triggered by stress, heat, or exercise
  \item Wearing closed-toe shoes for long periods, especially non-breathable materials like synthetic leather or plastic
  \item Wearing the same shoes every day without allowing them to dry out
  \item Poor foot hygiene — not washing or drying feet thoroughly
  \item Fungal infections such as athlete's foot (tinea pedis) — can worsen odor
  \item Hormonal changes (teenagers and pregnant women sweat more)
  \item Stress and anxiety — can trigger increased sweating
  \item Certain medications (some antidepressants and anticholinesterases)
\end{itemize}

\section*{When to See a Doctor}
See your doctor if:
\begin{itemize}
  \item Foot odor is severe and does not improve with good hygiene and over-the-counter treatments
  \item You have signs of a fungal infection: Itching, peeling, cracking, or redness between the toes
  \item You have signs of a bacterial infection: Redness, warmth, swelling, or pain in the foot
  \item You also have excessive sweating of the palms or underarms (may indicate general hyperhidrosis)
  \item The odor is accompanied by wounds or sores on your feet that are slow to heal (especially if you have diabetes)
\end{itemize}

\section*{Related Symptoms}
\begin{itemize}
  \item Excessive foot sweating
  \item Itching or burning between the toes
  \item Peeling or cracked skin on the feet
  \item Foot redness or tenderness
  \item General body odor or excessive sweating elsewhere
\end{itemize}
\textbf{Important:} Smelly feet are usually harmless, but if you have diabetes, even minor foot problems can become serious — check your feet daily and see a doctor promptly for any wounds, redness, or signs of infection.

\section*{Self-Care / Home Management}
\begin{itemize}
  \item Wash your feet daily with antibacterial soap, paying attention between the toes.
  \item Dry your feet thoroughly, especially between the toes, after washing.
  \item Change socks at least once a day — more often if your feet sweat heavily.
  \item Wear moisture-wicking socks made of natural fibers (cotton, wool) or synthetic athletic socks designed to pull moisture away from the skin.
  \item Alternate your shoes — do not wear the same pair two days in a row; allow them 24 hours to dry out.
  \item Use foot powder or antifungal powder to keep feet dry.
  \item Sprinkle baking soda or place cedar shoe inserts inside shoes to absorb odor.
  \item Soak your feet in black tea (the tannins reduce sweating) for 15-20 minutes daily for a week.
  \item Use over-the-counter antiperspirant sprays or roll-ons designed for feet.
  \item Wear open-toed shoes or sandals when possible to allow air circulation.
\end{itemize}

\section*{How Your Doctor Will Evaluate You}
\begin{itemize}
  \item Your doctor will examine your feet for signs of fungal infection, bacterial infection, or skin conditions.
  \item They will ask about your hygiene habits, shoe choices, and whether you have excessive sweating elsewhere.
  \item If a fungal infection is suspected, they may take a skin scraping to examine under a microscope.
  \item If general hyperhidrosis is suspected, they may recommend tests to rule out underlying conditions such as an overactive thyroid.
\end{itemize}

\section*{Treatment Options}
\begin{itemize}
  \item Most cases resolve with improved foot hygiene and proper footwear.
  \item Antifungal creams or sprays (clotrimazole, miconazole, terbinafine) if a fungal infection is present.
  \item Prescription-strength antiperspirants containing aluminum chloride hexahydrate for severe sweating.
  \item Iontophoresis — a device that passes a mild electrical current through water to reduce sweating; used for severe hyperhidrosis.
  \item Botox injections can temporarily block the nerves that stimulate sweat glands in the feet, providing relief for 6-12 months.
  \item Oral medications (anticholinergics) may be prescribed for severe, generalized hyperhidrosis.
  \item In very severe cases, surgical sympathectomy (cutting the nerves that trigger sweating) is an option, but it is rarely recommended.
\end{itemize}

\section*{References}
\begin{thebibliography}{99}
\bibitem{jaad-bromodosis} Moriarty B, et al. "Bromodosis: diagnosis and management." \emph{J Am Acad Dermatol}. 2024;91(1):112-120.
\bibitem{uptodate-foot-odor} Elewski BE, et al. Approach to the patient with foot odor. In: UpToDate, Post TW (Ed), UpToDate, Waltham, MA; 2025.
\bibitem{cochrane-foot-hyperhidrosis} Henteleff HJ, et al. "Interventions for plantar hyperhidrosis." \emph{Cochrane Database Syst Rev}. 2023;(5):CD013754.
\bibitem{jaad-hyperhidrosis} Pariser DM, et al. "Hyperhidrosis: clinical practice guideline." \emph{J Am Acad Dermatol}. 2023;88(5):1061-1070.
\bibitem{bmj-foot-care} Crawford F, et al. "Foot problems in adults." \emph{BMJ}. 2024;384:e077537.
\end{thebibliography}
\clearpage
